\chapter{1998年全国卷（文）}
\section{单选题}

\begin{problem}
\(\sin 600^{\circ}\) 的值是（\quad）

\choices
    {\(\frac{1}{2} \)}
    {\(-\frac{1}{2}\)}
    {\(\frac{\sqrt{3}}{2}\)}
    {\(-\frac{\sqrt{3}}{2}\)}

\end{problem}

\begin{problem}
函数 \(y=a^{|x|}\)（\(a>1\)）的图象是（\quad）

\end{problem}

\begin{problem}
已知直线 \(x=a\)（\(a>0\)）和圆 \((x - 1)^2 + y^2 = 4\) 相切，那么 \(a\) 的值是（\quad）

\choices
    {5}
    {4}
    {3}
    {2}

\end{problem}

\begin{problem}
两条直线 \(A_{1}x + B_{1}y + C_{1} = 0, A_{2}x + B_{2}y + C_{2} = 0\) 垂直的充要条件是（\quad）

\choices
    {\(A_{1}A_{2} + B_{1}B_{2} = 0\)}
    {\(A_{1}A_{2} - B_{1}B_{2} = 0\)}
    {\(\frac{A_1A_2}{B_1B_2} = -1\)}
    {\(\frac{B_1B_2}{A_1A_2} = 1\)}

\end{problem}

\begin{problem}
函数 \(f(x)=\frac{1}{x}\)（\(x\neq 0\)）的反函数 \(f^{-1}(x)=\)（\quad）

\choices
    {\(x\)（\(x\neq 0\)）}
    {\(\frac{1}{x}\)（\(x\neq 0\)）}
    {\(-x\)（\(x\neq 0\)）}
    {\(-\frac{1}{x}\)（\(x\neq 0\)）}

\end{problem}

\begin{problem}
已知点 \(P(\sin \alpha - \cos \alpha, \tan \alpha)\) 在第一象限，则在 \([0, 2\pi]\) 内 \(\alpha\) 的取值范围是（\quad）

\choices
    {\(\left(\frac{\pi}{2}, \frac{3\pi}{4}\right)\)}
    {\(\left(\pi, \frac{5\pi}{4}\right)\) \(\left(\frac{\pi}{4}, \frac{5\pi}{4}\right)\)}
    {\(\left(\frac{\pi}{2}, \frac{3\pi}{4}\right) \cup \left(\frac{5\pi}{4}, \frac{3\pi}{2}\right)\)}
    {\(\left(\frac{\pi}{4}, \frac{5\pi}{2}\right) \cup \left(\frac{3\pi}{4}, \frac{\pi}{2}\right)\)}

\end{problem}

\begin{problem}
已知圆锥体的全面积是底面积 \(3\) 的格点，则该圆锥的侧面展开图扇形的圆心角为（\quad）

\choices
    {\(120^{\circ}\)}
    {\(150^{\circ}\)}
    {\(180^{\circ}\)}
    {\(240^{\circ}\)}

\end{problem}

\begin{problem}
复数 \(\i\) 的一个立方根是 \(\i\)，它的另外两个立方根是（\quad）

\choices
    {\(\frac{\sqrt{3}}{2} \i + \frac{1}{2} \i\)}
    {\(-\frac{\sqrt{3}}{2} \i + \frac{1}{2} \i\)}
    {\(\frac{\sqrt{3}}{2} \i + \frac{1}{2} \i\)}
    {\(\pm \frac{\sqrt{3}}{2} \i - \frac{1}{2} \i\)}

\end{problem}

\begin{problem}
如果微分的两底面积分别是 \(S, S'\)，中截面的面积是 \(S_0\)，那么（\quad）

\choices
    {\(2\sqrt{S_0} = \sqrt{S} +\sqrt{S'}\)}
    {\(S_0 = \sqrt{SS'}\)}
    {\(2S_{0} = S + S^{\prime}\)}
    {\(S_0^2 = 2SS'\)}

\end{problem}

\begin{problem}
2名医生和4名护士被分配到2所学校为学生体检. 每校分配1名医生和2名护士，不同的分配方法共有（\quad）

\choices
    {6 种}
    {12 种}
    {18 种}
    {24 种}

\end{problem}

\begin{problem}
向高为 \(H\) 的水瓶中注水，注满为止，如果注水量 \(V\) 与水深 \(h\) 的函数关系的图象如下图所示，那么水瓶的形状是（\quad）

\end{problem}

\begin{problem}
椭圆 \(\frac{x^2}{12} + \frac{y^2}{3} = 1\) 的焦点为 \(F_1\) 和 \(F_2\). 点 \(P\) 在椭圆上，如果线段 \(PF_1\) 的中点 \(M\) 在 \(y\) 轴上，那么上点 \(M\) 的纵坐标是（\quad）

\choices
    {7 倍}
    {5 倍}
    {4 倍}
    {3 倍}

\end{problem}

\begin{problem}
球面上有 3 个点，其中任意两点的球面距离都等于大圆周长的 \(\frac{1}{6}\)，经过这 3 个点的小圆的周长为 \(4\pi\)，那么这个球的半径为（\quad）

\choices
    {\(\pm \frac{\sqrt{3}}{4}\)}
    {\(\pm \frac{\sqrt{3}}{2}\)}
    {\(\pm\frac{\sqrt{3}}{3} \)}
    {\(\pm\frac{3}{4} \)}

\end{problem}

\begin{problem}
一个直角三角形三内角的正弦值或等比数列，其最小内角的正弦值为（\quad）

\choices
    {\(\frac{1}{2} \)}
    {\(\frac{\sqrt{3}-1}{2} \)}
    {\(\frac{\sqrt{3}-2}{2} \)}
    {\(\frac{\sqrt{2\sqrt{3}+2}}{2} \)}

\end{problem}

\begin{problem}
等比数列 \(\{a_{n}\}\) 的公比为 \(-\frac{1}{2}\)，前 \(n\) 项和为 \(S_{n}\)，满足 \(\displaystyle\lim_{n\to \infty} S_n = \frac{1}{a_n}\)，那么 \(a_1\) 的值为（\quad）

\choices
    {\(\pm \sqrt{3}\)}
    {\(\pm \frac{3}{4}\)}
    {\(\pm\sqrt{2} \)}
    {\(\pm \frac{\sqrt{6}}{2}\)}

\end{problem}

\section{填空题}

\begin{problem}
设圆过双曲线 \(\frac{x^{2}}{a^{2}} - \frac{y^{2}}{4} = 1 \) 的一个顶点和一个焦点，圆心在此双曲线上，则圆心到双曲线中心的距离是 \fillinblank{}.

\end{problem}

\begin{problem}
\((x + 2)^{10}(x^{2} - 1)\) 的展开式中 \(x\) 的系数为 \fillinblank{}. （用数字作答）

\end{problem}

\begin{problem}
在直四棱柱 \(A_{1}B_{1}C_{1}D_{1} - ABCD\) 中，当底面四边形 \(ABCD\) 满足条件\fillinblank{}时，有 \(A_{1}C \perp B_{1}D_{1}\). （注：写上你认为正确的一种条件即可，不必考虑所有可能的情形）

\end{problem}

\begin{problem}
关于函数 \(f(x)=4\sin \left(2x+\frac{\pi}{3}\right)\)（\(x\in\R\)），有下列命题： \circled{1} 由 \(f(x_{1}) = f(x_{2}) = 0\) 可得 \(x_{1} - x_{2}\) 必是 \(\mathbf{r}\) 的整数倍； \circled{2} \(y = f(x)\) 的表达式可改写为 \(y = 4\cos \left(2x - \frac{\pi}{6}\right)\); \circled{3} \(y = f(x)\) 的图象交于点 \(\left(-\frac{\pi}{6},0\right)\) 对称； \circled{4} \(y = f(x) \) 的图象关于直线 \(x = -\frac{\pi}{6} \) 对称. 其中正确的命题的序号是 \fillinblank{}.（注：把你认为正确的命题的序号都填上）

\end{problem}

\begin{problem}
设 \(\alpha \neq b\)，解关于 \(x\) 的不等式: \(\alpha^2 x + b^2 (1 - x) \geqslant [\alpha x + b(1 - x)]^2\) \fillinblank{}.

\end{problem}

\begin{problem}
在 \(\triangle ABC\) 中，\(a, b, c\) 分别是 \(A, B, C\) 的对边，设 \(a + c = 2b\)，\(A - C = \frac{\pi}{3}\)，\(\frac{1}{3} \sin B\) 的值 \fillinblank{}.

\end{problem}

\section{解答题}

\begin{problem}
如图，直线 \(l_{1}\) 和 \(l_{2}\) 相交于点 \(M, l_{1} \perp l_{2}\)，点 \(N \in l_{1}\)，以 \(A, B\) 为端点的曲线段 \(C\) 上的任一点到 \(l_{2}\) 的距离与到点 \(N\) 的距离相等. 若 \(\triangle AMN\) 为锐角三角形，\(|AM| = \sqrt{17}, |AN| = 3\)，且 \(|BN| = 6\). 建立适当的坐标系，求曲线段 \(C\) 的方程.

\end{problem}

\begin{problem}
已知斜三棱柱 \(ABC - A_{1}B_{1}C_{1}\) 的侧面 \(A_{1}ACC_{1}\) 与底面 \(ABC\) 垂直，\(\angle ABC = 90^{\circ}\)，\(BC = 2\)，\(AC = 2\sqrt{3}\). 且 \(A_{1}A_{1} \perp A_{1}C_{1}\)，\(AA_{1} = A_{1}C\). （1）求侧面 \(A_{1}A \) 与底面 \(ABC\) 所成角的大小; （2）求侧面 \(A_{1}ABB_{1} \) 与底面 \(ABC\) 所成二面角的大小; （3）求侧棱 \(B_{1}B \) 和侧面 \(A_{1}ACC_{1} \) 的距离.

\end{problem}

\begin{problem}
如图，为处理含有某种杂质的污水，要制造一底宽为 2 米的无盖长方体沉淀槽，污水从 \(A\) 点流入，经沉淀后从 \(B\) 点流出. 设桶体的长度为 \(a\) 米，高度为 \(b\) 米. 已知流出的水中该杂质的质量分数与 \(a\)，\(b\) 的乘积 \(ab\) 成反比. 现有制桶材料 60 平方米. 当问 \(a\)，\(b\) 各为多少米时，经沉淀后流出的水中该杂质的质量分数最小. （\(A\)、\(B\) 孔的面积忽略不计）
\end{problem}

